\documentclass[pdflatex,sn-mathphys-ay]{sn-jnl}

\usepackage{graphicx}
\usepackage{booktabs}
\usepackage{amsmath,amssymb,amsfonts}
\usepackage{multirow}

\jyear{2026}

\raggedbottom

\begin{document}

\title[FALI]{FALI: Mechanistic Decomposition of Conditional Attention-Head Importance}

% Replace all author metadata before submission.
\author*[1]{\fnm{Author} \sur{Name}}\email{author@example.com}
\affil*[1]{\orgdiv{Department}, \orgname{Institution}, \orgaddress{\city{City}, \country{Country}}}

\abstract{%
% TARGET: 150--250 words. Draft last, after all claims/results are frozen.
% Required content: problem; gap; FALI estimand and factorization; validation;
% GPT-2 IOI result; Pythia induction reversal; bounded conclusion.
}

\keywords{mechanistic interpretability, causal intervention, transformer circuits, attention heads, conditional importance}

\maketitle

\section{Introduction}\label{sec:introduction}
% Target: ~1.75--2 pages.
% Arc: conditional importance is known -> mechanism remains unresolved -> FALI ->
% validation -> GPT-2/Pythia mechanism reversal -> contributions.

\section{Related Work}\label{sec:related}

\subsection{Self-repair and conditional component importance}\label{sec:rw-conditional}
% Hydra, copy suppression, self-repair, CoAx.

\subsection{Interaction-aware causal attribution}\label{sec:rw-interaction}
% Activation/path patching, CHG, Curse of Multiple Mediators.

\subsection{Decomposing attention computation}\label{sec:rw-attention}
% QK/OV, causal attention communication, QK/value decompositions, Owen machinery.

\section{Problem Formulation}\label{sec:problem}

\subsection{Conditional causal usefulness}\label{sec:conditional-usefulness}
% Define intact state A, perturbed state B, candidate head h, output-grounded
% usefulness, intact effect, conditional effect, conditional gain.

\subsection{Mechanism-resolution objective}\label{sec:mechanism-objective}
% State that FALI assumes the candidate head/target edge are already selected.

\section{Factorized Attention-Level Intervention}\label{sec:fali}

\subsection{Architecture-aligned factors}\label{sec:factors}
% T target score(s); R rival scores; O other structurally valid scores;
% V value content; C surrounding-network/context.

\subsection{Counterfactual intervention game}\label{sec:game}
% Define U(T,R,O,V,C) as output-grounded candidate usefulness in the same C context.
% State explicit off-manifold/counterfactual semantics.

\subsection{Standard Owen attribution and coarse--fine reconciliation}\label{sec:owen}
% Partition: SCORE={T,R,O}, VALUE={V}, CONTEXT={C}.
% State efficiency and quotient identities. No mathematical novelty claim.

\section{Method Validation}\label{sec:validation}

\subsection{Analytical and synthetic ground truth}\label{sec:synthetic}
% Five single-channel cases and cross-factor interaction tests.

\subsection{Transformer implementation invariants}\label{sec:implementation-validation}
% Batch-safe masks, endpoint identity, score->pattern, z reconstruction,
% no-live-candidate leak, untouched-head checks, quotient bridge.

\section{GPT-2 IOI: Conditional Recruitment and Fine-Score Decomposition}\label{sec:gpt2}

\subsection{Experimental design}\label{sec:gpt2-design}
% Canonical primary heads; candidate heads; 15 templates x BABA/ABBA;
% four seeds; output metric; preregistered confirmation.

\subsection{Conditional recruitment is not attention increase}\label{sec:gpt2-functional}
% Corrected functional baseline; L10H2 genuine recruitment; L10H10 contrast.

\subsection{Canonical FALI mechanism}\label{sec:gpt2-fali}
% L10H2 and L10H10: T positive, R negative; O descriptive;
% coarse score/value/context only where necessary.

\subsection{Intervention robustness}\label{sec:gpt2-robustness}
% Report matched-mean and matched-resample fine-channel robustness accurately.
% Explicitly disclose that the preregistered overall EF5 gate failed because the
% zero-damage reference criterion was not satisfied.

\section{Cross-Setting Test: Pythia Induction}\label{sec:pythia}

\subsection{Preregistered transfer design}\label{sec:pythia-design}
% Pythia-160M induction; primary L5H0; candidate L8H2; frozen transfer hypotheses.

\subsection{Recruitment generalizes but the fine mechanism reverses}\label{sec:pythia-results}
% conditional effect/gain positive; phi_T negative; phi_R positive; transfer gate failed.

\section{Discussion}\label{sec:discussion}

\subsection{A common functional response need not imply a common mechanism}\label{sec:discussion-mechanisms}

\subsection{Relationship to conditional attribution and patching}\label{sec:discussion-positioning}
% Localization/CoAx: which; interaction methods: conditionality; FALI: mechanism resolution.

\section{Limitations}\label{sec:limitations}
% Required: hypothesis-driven target selection; off-manifold hybrids; coarse C;
% partition/intervention dependence; only two substantive model/task settings.

\section{Conclusion}\label{sec:conclusion}
% Compact: what FALI adds and what the cross-setting reversal establishes.

\section*{Statements and Declarations}

\subsection*{Competing Interests}
% TODO: author-provided statement.

\subsection*{Funding}
% TODO: author-provided statement.

\subsection*{Data Availability}
% TODO: state exact public/private data availability at submission.

\subsection*{Code Availability}
% TODO: repository/archive and reproducibility statement.

\subsection*{Author Contributions}
% TODO: author-provided CRediT-style contribution statement if appropriate.

\subsection*{Use of Generative AI}
% TODO before submission: document generative-AI use to the extent required by
% current Springer/Machine Learning policy and the final actual workflow.

\bibliography{references}

\end{document}
